\documentclass[11pt]{article}
\usepackage[hmargin=0.85in,vmargin=0.8in]{geometry}
\usepackage{amsmath,amssymb,amsthm}
\usepackage[T1]{fontenc}
\usepackage{lmodern}
\usepackage[hidelinks]{hyperref}
\setlength{\emergencystretch}{1em}
\newtheorem{lemma}{Lemma}
\newtheorem{theorem}{Theorem}
\newcommand{\Q}{\mathbb Q}
\newcommand{\C}{\mathbb C}
\newcommand{\Hh}{\mathcal H}
\newcommand{\SL}{\mathrm{SL}_2(\mathbb Z)}
\title{AI-assisted disproof of conjecture 00000009679}
\author{idealistichacker}
\date{October 29, 2026}
\begin{document}
\maketitle
\begin{abstract}
With the standard complex-vector-space meaning of cusp form, the first
universal clause of conjecture 00000009679 fails. The algebraic point
$\tau=\sqrt2+i$ is outside every quadratic subfield of $\C$. The actual
level-one weight-12 discriminant $\Delta=\eta^{24}$ is nonzero at $\tau$;
thus the nonzero cusp form $f=\Delta(\tau)^{-1}\Delta$ has the algebraic
value $f(\tau)=1$. This is evaluation rescaling, not Fourier or Hecke
normalization. No algebraicity of the coefficients of $f$ is asserted.
We separate the proof and recorded Lean core checks from full portable
package reproduction. Independent PDF and final review remain necessary.
\end{abstract}
\section{The Clause and Its Hypotheses}
Both source languages assert that every nonzero cusp modular form of
weight $k\geq2$ has transcendental values at algebraic points outside
imaginary quadratic fields. We use the standard analytic definition:
$S_k(\SL)$ is a complex vector space of holomorphic modular forms
vanishing at the cusps \cite[Definition 4.5, p.~49]{Milne}.
In particular, for all $\left(\begin{smallmatrix}a&b\\c&d\end{smallmatrix}\right)
\in\SL$, a member satisfies
$f((az+b)/(cz+d))=(cz+d)^k f(z)$ on $\Hh$.
Level one is already a valid instance of the asserted unrestricted class.
Neither language imposes algebraic Fourier coefficients, a fixed
coefficient field, first Fourier coefficient one, or a Hecke normalization.

Writing $\overline{\Q}$ for the algebraic numbers in $\C$, the necessary
level-one clause is
\[
 \begin{gathered}
 \forall k\geq2\ \forall f\in S_k(\SL)\setminus\{0\}\
 \forall z\in\Hh\cap\overline{\Q},\\
 \bigl(z\text{ lies in no imaginary quadratic subfield of }\C\bigr)
 \Longrightarrow f(z)\notin\overline{\Q}.
 \end{gathered}
\tag{U}
\]
We will fix one pair $(f,\tau)$, not vary the form during evaluation.
The later CM-period/class-number clause is not analyzed; falsifying the
necessary clause (U) suffices for the combined assertion as written.
\section{An Algebraic Point Outside Every Quadratic Field}
\begin{lemma}
For the positive real square root $r=\sqrt2$, put $\tau=r+i$.
Then $\tau\in\Hh\cap\overline{\Q}$, and $\tau$ belongs to no
 degree-two intermediate field of $\C/\Q$.
\end{lemma}
\begin{proof}
We have $\operatorname{Im}\tau=1>0$ and
\[
 \tau^2=1+2ri,\qquad (\tau^2-1)^2=-8,
 \qquad \tau^4-2\tau^2+9=0.
\]
The last identity is a nonzero rational polynomial relation, proving
algebraicity. Suppose $a\tau^2+b\tau+c=0$ with $a,b,c\in\Q$.
Its imaginary part is $2ar+b=0$. If $a\ne0$, then
$r=-b/(2a)\in\Q$, contradicting the irrationality of $\sqrt2$.
Thus $a=0$, then $b=0$, and finally $c=0$. Consequently
$1,\tau,\tau^2$ are linearly independent over $\Q$.
If $\tau\in K\subseteq\C$ and $[K:\Q]=2$, all three elements lie in
 the two-dimensional rational vector space $K$, a contradiction.
This excludes every such $K$, in particular every imaginary quadratic
field; excluding only $\Q(i)$ would not suffice.
\end{proof}
\newpage
\section{The Actual Cusp Form and the Contradiction}
For $z\in\Hh$, put $q=e^{2\pi iz}$ and use the convention
\[
 \Delta(z)=\eta(z)^{24}=q\prod_{n=1}^{\infty}(1-q^n)^{24}.
\]
This is the classical level-one weight-12 cusp form, with first Fourier
coefficient one, and it has no zeros in $\Hh$.
Milne's discriminant has a different nonzero constant factor.
His \S5 (p.~67) states the weight-12 cusp-form property and identifies
it with $(2\pi)^{12}\eta^{24}$; Proposition 3.12 (p.~47) establishes
nonvanishing of the lattice discriminant \cite{Milne}.
Dividing by that constant gives our convention. Any nonzero constant
factor cancels in the evaluation rescaling below.

Nonvanishing of the product is not inferred merely from nonvanishing
of its factors. Since $0<|q|<1$, the sum $\sum_{n\geq1}|q|^n$ converges.
For sufficiently large $n$, the analytic logarithm of $1-q^n$ near one
has absolute value at most $2|q|^n$. The sum of these logarithms converges
absolutely, so its exponential is nonzero. The finitely many omitted
factors and $q$ are also nonzero. Hence $\Delta(z)\ne0$.
\begin{theorem}
There exist a nonzero $f\in S_{12}(\SL)$ and an algebraic
$\tau\in\Hh$ outside every imaginary quadratic subfield of $\C$ such
that $f(\tau)=1$. In particular, (U) is false.
\end{theorem}
\begin{proof}
Choose the point in Lemma 1 and set
\[
 v=\Delta(\tau)\ne0,\qquad c=v^{-1},\qquad f(z)=c\Delta(z).
\]
A fixed complex scalar preserves the same weight-12 transformation law,
holomorphy on $\Hh$, and vanishing at all cusps.
Thus $f\in S_{12}(\SL)$, and $f(\tau)=v^{-1}v=1$; in particular $f\ne0$.
Lemma 1 gives every required condition on $\tau$, and $12\geq2$.
But $1\in\Q$ is algebraic, not transcendental over $\Q$, contradicting (U).
\end{proof}
\paragraph{Hypothesis audit.}
The object is a genuine cusp form for $\SL$, not a constant function
or a mock evaluation map. Its weight is the integer 12. Its nonzero
status follows from its value 1. The point lies in the open upper
half-plane, is algebraic over $\Q$, and lies in no degree-two subfield.
No excluded CM point, zero form, boundary point, or vanishing
discriminant is used.
\paragraph{What this does not prove.}
Choosing $\tau$ first and then fixing $f$ once and for all is legitimate
when negating a universal assertion about both.
The first Fourier coefficient of $f$ is $c$, not an asserted 1, and
we prove neither that $c$ is algebraic nor that it is transcendental.
We do not claim algebraic Fourier coefficients or a Fourier-normalized
or Hecke-normalized arithmetic counterexample. Evaluation rescaling
must not be confused with those normalizations. This argument does not
overturn any theorem of Nesterenko, establish an arithmetic-restricted
variant, or settle a CM-period or class-number claim. An applicable
additional arithmetic convention would require a different analysis.
\newpage
\section{Lean Objects and the Evidence Boundary}
\raggedright
\small
All names below have namespace \texttt{ModularValue9679}. This is an
object map of the two complete sources read on October 29, 2026.
Recorded core replays below are not a new Lean run by this paper author.
\paragraph{Point module.}
\texttt{AlgebraicPoint.lean} defines \texttt{point} as $\sqrt2+i$.
Theorems \texttt{point\_im\_pos} and \texttt{point\_algebraic} give
positivity and algebraicity. The latter uses closure of algebraic numbers
under powers and addition, not the displayed quartic identity.
The theorems \texttt{quadratic\_relation\_trivial} and
\texttt{point\_triple\_independent} implement the independence argument.
Crucially, \texttt{point\_not\_mem\_quadratic} takes an actual
\texttt{IntermediateField Rat Complex} of rational \texttt{finrank} 2
and excludes membership by the dimension bound; non-CM status is not
assumed or encoded as an unexplained Boolean.
\paragraph{Cusp-form module.}
\texttt{NormalizedCusp.lean} defines \texttt{tau : UpperHalfPlane}
using positivity. Its \texttt{ImaginaryQuadraticSubfield} requires an
actual degree-two intermediate field and a nonreal element;
\texttt{OutsideImaginaryQuadratic} quantifies over all such fields.
The stronger point theorem supplies \texttt{tau\_outside\_imaginary\_quadratic}.
The object \texttt{normalized} is a complex scalar multiple of the actual
\texttt{CuspForm.discriminant} at level one and weight 12.
Mathlib's \texttt{ModularForm.discriminant\_ne\_zero}, with the actual
cusp-form coercion, supplies nonvanishing. The source gives
\texttt{normalized\_value}, \texttt{normalized\_nonzero}, and
\texttt{normalized\_value\_algebraic}. Finally,
\texttt{level\_one\_value\_assertion\_false} negates the explicit
universal clause, while \texttt{actual\_counterexample} packages the
form, point, and their properties.
\paragraph{Actual evidence and remaining work.}
The parent receipt dated 2026-10-27, 05:30:59 UTC records both modules
exiting zero in 33.813 seconds, with all 19 public theorem types and axiom
closures checked. A separate core review records an independent replay
on that date, 05:39:44--05:40:24 UTC, exiting zero in 41.032 seconds;
its logs include those 19 theorems and six definition types, bodies, and
axiom closures. All 25 recorded closures contain only \texttt{propext},
\texttt{Classical.choice}, and \texttt{Quot.sound}. This author checked
the actual receipts, logs, and matching source hashes, not a new Lean run.
The review result lists a report and two further evidence files that are
absent; none is relied on. These runs use Lean 4.33.1 and Mathlib
\texttt{\small 0df444a360eaa60ab8c11dca51a86af692955474}.
Shared pinned caches remain a trust boundary. Complete portable
reproduction with all nine physical dependencies, a final package,
independent PDF review, and publication clearance are not established by
this paper. Separate author receipts record actual Tectonic compilation
and page rendering, not an independent final review or upstream approval.
\paragraph{Source anchor and disclosure.}
The complete bilingual source is \texttt{conjectures/00000009679.md} at
this locally anchored version, not a fresh upstream or duplicate check:
\begin{quote}\footnotesize\ttfamily
Commit: b355e42906e6eaf52bc7e9d005e0ef02d8aa7400\\
Git blob: 541b54697f4283123781a066f17e9b0a7e3ec223\\
Raw SHA-256:\\
e3fd65548a676ba5dcee0340d2f9dd1aae6c2a069273e9d61799e4bba712b639
\end{quote}
This is an AI-assisted paper by the stated author. No affiliation,
independent final-review approval, claim, or publication is asserted.
\begin{thebibliography}{9}
\bibitem{Milne} J. S. Milne, \emph{Modular Functions and Modular Forms},
v1.31, March 22, 2017. Author's course notes, compact edition:
Definition 4.5, p.~49; Proposition 3.12, p.~47; product formula at the
start of \S5, p.~67. \url{https://www.jmilne.org/math/CourseNotes/MFc.pdf}.
\end{thebibliography}
\end{document}
